\ifdefined\maindoc\else
% typesetting this chapter as a standalone document
\def\doctitle{ArtiSynth Overview}
% starting definitions for both the main document and stand-alone chapters
\documentclass{book}

\def\mech{artisynth.core.mechmodels}
\def\mgeo{maspack.geometry}

% Add search paths for input files
\makeatletter
\def\input@path{{../}{../../}{../texinputs/}}
\makeatother

\usepackage{amsmath}
\usepackage{framed}
%%
%% Default settings for artisynth
%%
\NeedsTeXFormat{LaTeX2e}
%%\ProvidesPackage{artisynthDoc}[2012/04/05]

\usepackage[T1]{fontenc}
\usepackage[latin1]{inputenc}
\usepackage{listings}
\usepackage{makeidx}
\usepackage{latexml}
\usepackage{graphicx}
\usepackage{framed}
\usepackage{booktabs}
\usepackage{color}

\newcommand{\pubdate}{\today}
\newcommand{\setpubdate}[1]{\renewcommand{\pubdate}{#1}}
\newcommand{\code}[1]{{\tt #1}}

\iflatexml
\usepackage{hyperref}
\setlength\parindent{0pt} 
\else
%% then we are making a PDF, so include things that LaTeXML can't handle: 
%% docbook style, \RaggedRight
\usepackage{ifxetex}
\usepackage{xstring}
\usepackage{pslatex} % fixes fonts; in particular sets a better-fitting \tt font

\usepackage[most]{tcolorbox}
\definecolor{shadecolor}{rgb}{0.95,0.95,0.95}
\tcbset{
    frame code={}
    center title,
    left=0pt,
    right=0pt,
    top=0pt,
    bottom=0pt,
    colback=shadecolor,
    colframe=white,
    width=\dimexpr\textwidth\relax,
    enlarge left by=0mm,
    boxsep=0pt,
    arc=0pt,outer arc=0pt,
}%

\usepackage[A4]{artisynth_papersize}
%\usepackage[letter]{artisynth_papersize}
\usepackage[hyperlink]{asciidoc-dblatex} 

%\usepackage{verbatim}
\usepackage{ragged2e}
\setlength{\RaggedRightRightskip}{0pt plus 4em}
\RaggedRight
\renewcommand{\DBKpubdate}{\pubdate}
\renewcommand{\DBKreleaseinfo}{}
\fi

% set hypertext links to be dark blue:
\definecolor{darkblue}{rgb}{0,0,0.8}
\definecolor{sidebar}{rgb}{0.5,0.5,0.7}
\hypersetup{colorlinks=true,urlcolor=darkblue,linkcolor=darkblue,breaklinks=true}

%%%%%%%%%%%%%%%%%%%%%%%%%%%%%%%%%%%%%%%%%%%%%%%%%%%%%%%%%%%%%%%%%%%%%%%%%%%%%
%
% Define macros for handling javadoc class and method references
%
%%%%%%%%%%%%%%%%%%%%%%%%%%%%%%%%%%%%%%%%%%%%%%%%%%%%%%%%%%%%%%%%%%%%%%%%%%%%%
\makeatletter

% macro to enable line break if inside a PDF file
\def\pdfbreak{\iflatexml\else\\\fi}

% code inspired by http://stackoverflow.com/questions/2457780/latex-apply-an-operation-to-every-character-in-a-string
\def\removeargs #1{\doremoveargs#1$\wholeString\unskip}
\def\doremoveargs#1#2\wholeString{\if#1$%
\else\if#1({()}\else{#1}\taketherest#2\fi\fi}
\def\taketherest#1\fi
{\fi \doremoveargs#1\wholeString}

% Note: still doesn't work properly when called on macro output ...
% i.e., \dottoslash{\concatnames{model}{base}{foo}} fails 
\def\dottoslash #1{\dodottoslash#1$\wholeString\unskip}
\def\dodottoslash#1#2\wholeString{\if#1$%
\else\if#1.{/}\else{#1}\fi\dottaketherest#2\fi}
\def\dottaketherest#1\fi{\fi \dodottoslash#1\wholeString}

\def\hashtodot #1{\dohashtodot#1$\wholeString\unskip}
\def\dohashtodot#1#2\wholeString{\if#1$X%
\else\if#1\#{.}\else{#1}\fi\hashtaketherest#2\fi}
\def\hashtaketherest#1\fi{\fi \dohashtodot#1\wholeString}

%\dollartodot{#1} does the same thing as \StrSubstitute[0]{#1}{\$}{.}
% from the packahe xstring. We define \dollartodot instead because
% LaTeXML does not implement xstring.
%
% Note that for the substituion to work, we need \ifx instead of \if,
% since otherwise escaped characters won't work properly:
% if #1 = \$, then \if#1* seems to compare '\' and '$' (and output '*'),
% rather than comparing '$' to '*'
\def\dollartodot #1{\dodollartodot#1*\wholeString\unskip}
\def\dodollartodot#1#2\wholeString{\ifx#1*%
\else \ifx#1\${.}\else{#1}\fi\dollartaketherest#2\fi}
\def\dollartaketherest#1\fi{\fi \dodollartodot#1\wholeString}

% concatenates up to three class/method names together, adding '.' characters
% between them. The first and/or second argument may be empty, in which case
% the '.' is omitted. To check to see if these arguments are empty, we
% use a contruction '\if#1@@', which will return true iff #1 is empty
% (on the assumption that #1 will not contain a '@' character).
\def\concatnames
#1#2#3{\if#1@@\if#2@@#3\else #2.#3\fi\else\if#2@@#1.#3\else#1.#2.#3\fi\fi}

\newcommand{\javabase}{}
\newcommand{\setjavabase}[1]{\renewcommand{\javabase}{#1}}

\def\artisynthDocBase{@ARTISYNTHDOCBASE}

\iflatexml
\def\ifempty#1{\def\temp{#1}\ifx\temp\empty}%
\newcommand{\artisynthManual}[3][]{%
   \ifempty{#1}
      \href{@ARTISYNTHDOCBASE/#2/#2.html}{#3}%
    \else
      \href{@ARTISYNTHDOCBASE/#1/#2.html}{#3}%
    \fi
}
\else
\newcommand{\artisynthManual}[3][]{%
\href{https://www.artisynth.org/@ARTISYNTHDOCBASE/#2.pdf}{#3}}
\fi

%\href{@ARTISYNTHDOCBASE/#2/#2.html}{#3}}



\newcommand{\javaclassx}[2][]{%
% Includes code to prevent an extra '.' at the front if #1 is empty. It
% works like this: if '#1' is empty, then '#1.' expands to '.', and so 
% '\if#1..' will return true, in which case we just output '#2'.
\href{@JDOCBEGIN/\concatnames{\javabase}{#1}{#2}@JDOCEND}{#2}}
\newcommand{\javaclass}[2][]{%
\href{@JDOCBEGIN/\concatnames{}{#1}{#2}@JDOCEND}{\dollartodot{#2}}}
\newcommand{\javaclassAlt}[2]{%
\href{@JDOCBEGIN/\concatnames{}{}{#1}@JDOCEND}{#2}}

\newcommand{\javamethodArgsx}[2][]{%
\href{@JDOCBEGIN/\concatnames{\javabase}{#1}{#2}@JDOCEND}{#2}}
\newcommand{\javamethodArgs}[2][]{%
\href{@JDOCBEGIN/\concatnames{}{#1}{#2}@JDOCEND}{#2}}
\newcommand{\javamethodAlt}[2]{%
\href{@JDOCBEGIN/\concatnames{}{}{#1}@JDOCEND}{#2}}
\newcommand{\javamethodAltx}[2]{%
\href{@JDOCBEGIN/\concatnames{\javabase}{}{#1}@JDOCEND}{#2}}

\newcommand{\javamethodNoArgsx}[2][]{%
\href{@JDOCBEGIN/\concatnames{\javabase}{#1}{#2}@JDOCEND}{\removeargs{#2}}}
\newcommand{\javamethodNoArgs}[2][]{%
\href{@JDOCBEGIN/\concatnames{}{#1}{#2}@JDOCEND}{\removeargs{#2}}}

\newcommand{\javamethod}{\@ifstar\javamethodNoArgs\javamethodArgs}
\newcommand{\javamethodx}{\@ifstar\javamethodNoArgsx\javamethodArgsx}

%%%%%%%%%%%%%%%%%%%%%%%%%%%%%%%%%%%%%%%%%%%%%%%%%%%%%%%%%%%%%%%%%%%%%%%%%%%%%
%
% Define macros for sidebars
%
%%%%%%%%%%%%%%%%%%%%%%%%%%%%%%%%%%%%%%%%%%%%%%%%%%%%%%%%%%%%%%%%%%%%%%%%%%%%%

\iflatexml
\newenvironment{sideblock}{\begin{quote}}{\end{quote}}
\else
\usepackage[strict]{changepage}
\definecolor{sidebarshade}{rgb}{1.0,0.97,0.8}
\newenvironment{sideblock}{%
    \def\FrameCommand{%
    \hspace{1pt}%
    {\color{sidebar}\vrule width 2pt}%
    %{\vrule width 2pt}%
    {\color{sidebarshade}\vrule width 4pt}%
    \colorbox{sidebarshade}%
  }%
  \MakeFramed{\advance\hsize-\width\FrameRestore}%
  \noindent\hspace{-4.55pt}% disable indenting first paragraph
  \begin{adjustwidth}{}{7pt}%
  %\vspace{2pt}\vspace{2pt}%
}
{%
  \vspace{2pt}\end{adjustwidth}\endMakeFramed%
}
\fi

\iflatexml
\newenvironment{shadedregion}{%
  \definecolor{shadecolor}{rgb}{0.96,0.96,0.98}%
  \begin{shaded*}%
% Put text inside a quote to create a surrounding blockquote that
% will properly accept the color and padding attributes
  \begin{quote}%
}
{%
  \end{quote}%
  \end{shaded*}%
}
\else
\newenvironment{shadedregion}{%
  \definecolor{shadecolor}{rgb}{0.96,0.96,0.98}%
  \begin{shaded*}%
}
{%
  \end{shaded*}%
}
\fi

% Wanted to create a 'listing' environment because lstlisting is
% tedious to type and because under latexml it may need
% some massaging to get it to work properly. But hard to do
% because of the verbatim nature of listing
%\iflatexml
%\newenvironment{listing}{\begin{lstlisting}}{\end{lstlisting}}%
%\else
%\newenvironment{listing}{\begin{lstlisting}}{\end{lstlisting}}%
%\fi

\iflatexml\else
% fancyhdr was complaining that it wanted a 36pt header height ...
\setlength{\headheight}{36pt}
\fi

% macro for backslash character
\newcommand\BKS{\textbackslash}

% macro for double hyphen (to prevent conversion of -- into -)
\newcommand\DHY{-{}-}

% Convenience stuff
\newcommand{\ifLaTeXMLelse}[2]{%
  \iflatexml %
  #1 %
  \else %
  #2 %
  \fi %
}

\newcommand{\ifLaTeXML}[1]{ %
  \iflatexml %
  #1 %
  \fi %
}

% new methodtable environment for documenting methods

% base width of the method table
\newlength{\methodtablewidth}
\iflatexml
\setlength{\methodtablewidth}{1.4\textwidth}
\else
\setlength{\methodtablewidth}{0.94\textwidth}
\fi
% horizontal space added at end of call to \methodentry
\newlength{\methodskip}
\setlength{\methodskip}{0pt}
% lengths set inside methodtable environment:
\newlength{\methodsiglength} % length of the method signature
\newlength{\methodcomlength} % length of the method comment
\setlength{\methodsiglength}{0.5\methodtablewidth}
\setlength{\methodcomlength}{0.5\methodtablewidth}

% command to add a method to a method table:
% arg #1: package and signature for finding URL
% arg #2: anchor text
% arg #3: comment describing the method
\newcommand{\methodentry}[3]{%
\javamethodAlt{#1}{\parbox[t]{\methodsiglength}{#2}}&
{\parbox[t]{\methodcomlength}{#3}}\\%
\noalign{\vspace{\methodskip}}}

% methodtable environment takes two arguments, both scale factors for
% methodtablewidth:
% arg #1: width of the method signature column
% arg #2: width of the method comment column
\newenvironment{methodtable}[3][0pt]{%
\begingroup
\setlength{\topskip}{0pt}
\setlength{\methodskip}{#1}
\setlength{\methodsiglength}{#2\methodtablewidth}%
\setlength{\methodcomlength}{#3\methodtablewidth}%
\iflatexml
\begin{snugshade}
\else
\begin{tcolorbox}
\fi
\renewcommand{\arraystretch}{1}
\begin{tabular}{ll}}{%
\end{tabular}
\renewcommand{\arraystretch}{1}
\iflatexml
\end{snugshade}
\else
\end{tcolorbox}
\fi
\endgroup}

% commands for added top, mid and bottom lines in the table.
% uses booktabs for PDF, regular hline for HTML
\newcommand{\topline}{\iflatexml\hline\else\toprule\fi}
\newcommand{\midline}{\iflatexml\hline\else\midrule\fi}
\newcommand{\botline}{\iflatexml\hline\else\bottomrule\fi}
\newcommand{\blankline}{%
\multicolumn{2}{l}{\iflatexml{@SPACE}\else\phantom{M}\fi}\\}%
% add vertical space within a two colum method environment
\newcommand{\methodspace}[1]{%
\iflatexml
\multicolumn{2}{l}{@VERTSPACE[#1]}\\
\else
\noalign{\vspace{#1}}%
\fi}%
% break a line and add an indentation of 1em
\newcommand{\brh}{\\\phantom{M}}

\makeatother

\def\matl{\left(\begin{matrix}}
\def\matr{\end{matrix}\right)}

\def\Bthe{\boldsymbol\theta}
\def\Btau{\boldsymbol\tau}
\def\Bom{\boldsymbol\omega}
\def\Bdel{\boldsymbol\delta}
\def\Blam{\boldsymbol\lambda}
\def\Bphi{\boldsymbol\phi}
\def\Bxi{\boldsymbol\xi}
\def\Bgam{\boldsymbol\gamma}
\def\Bsig{\boldsymbol\sigma}
\def\Bnu{\boldsymbol\nu}
\def\Bmu{\boldsymbol\mu}

\def\A{{\bf A}}
\def\B{{\bf B}}
\def\C{{\bf C}}
\def\D{{\bf D}}
\def\F{{\bf F}}
\def\G{{\bf G}}
\def\H{{\bf H}}
\def\I{{\bf I}}
\def\J{{\bf J}}
\def\K{{\bf K}}
\def\Jc{{\bf J}_c}
\def\L{{\bf L}}
\def\M{{\bf M}}
\def\N{{\bf N}}
\def\O{{\bf O}}
\def\P{{\bf P}}
\def\Q{{\bf Q}}
\def\R{{\bf R}}
\def\T{{\bf T}}
\def\U{{\bf U}}
\def\W{{\bf W}}
\def\X{{\bf X}}
\def\Minv{{\bf M}^{-1}}

\def\a{{\bf a}}
\def\b{{\bf b}}
\def\c{{\bf c}}
\def\d{{\bf d}}
\def\e{{\bf e}}
\def\f{{\bf f}}
\def\g{{\bf g}}
\def\k{{\bf k}}
\def\l{{\bf l}}
\def\m{{\bf m}}
\def\n{{\bf n}}
\def\p{{\bf p}}
\def\q{{\bf q}}
\def\r{{\bf r}}
\def\u{{\bf u}}
\def\v{{\bf v}}
\def\w{{\bf w}}
\def\x{{\bf x}}
\def\y{{\bf y}}
\def\z{{\bf z}}

\def\ma{{\bf m}_\alpha}
\def\mb{{\bf m}_\beta}
\def\va{{\bf v}_\alpha}
\def\vb{{\bf v}_\beta}
\def\vp{{\bf v}_\rho}
\def\vk{{\bf v}_k}
\def\ua{{\bf u}_\alpha}
\def\ub{{\bf u}_\beta}
\def\uk{{\bf u}_k}
\def\uj{{\bf u}_j}
\def\mar{{\bf m}_{\alpha r}}
\def\mbr{{\bf m}_{\beta r}}

\def\Maa{{\bf M}_{\alpha\alpha}}
\def\Mab{{\bf M}_{\alpha\beta}}
\def\Mba{{\bf M}_{\beta\alpha}}
\def\Mbb{{\bf M}_{\beta\beta}}
\def\hatMaa{\hat{\bf M}_{\alpha\alpha}}
\def\hatMab{\hat{\bf M}_{\alpha\beta}}
\def\hatMba{\hat{\bf M}_{\beta\alpha}}
\def\hatMbb{\hat{\bf M}_{\beta\beta}}
\def\Mbp{{\bf M}_{\beta\rho}}
\def\Map{{\bf M}_{\alpha\rho}}
\def\Mpa{{\bf M}_{\rho\alpha}}
\def\Mpb{{\bf M}_{\rho\beta}}
\def\Mpp{{\bf M}_{\rho\rho}}
\def\Mbk{{\bf M}_{\beta k}}
\def\Mak{{\bf M}_{\alpha k}}
\def\Mka{{\bf M}_{k\alpha}}
\def\Mkb{{\bf M}_{k\beta}}
\def\Mkk{{\bf M}_{kk}}

\def\Ga{{\bf G}_{\alpha}}
\def\Gp{{\bf G}_{\rho}}
\def\Gaa{{\bf G}_{\alpha\alpha}}
\def\Gab{{\bf G}_{\alpha\beta}}
\def\Gba{{\bf G}_{\beta\alpha}}
\def\Gbb{{\bf G}_{\beta\beta}}
\def\Gap{{\bf G}_{\alpha\rho}}
\def\Gpa{{\bf G}_{\rho\alpha}}
\def\Gbp{{\bf G}_{\beta\rho}}
\def\Gak{{\bf G}_{\alpha k}}
\def\Gka{{\bf G}_{k\alpha}}
\def\Gja{{\bf G}_{j\alpha}}
\def\Gkb{{\bf G}_{k\beta}}
\def\Gbk{{\bf G}_{\beta k}}

\def\lama{\Blam_{\alpha}}
\def\lamb{\Blam_{\beta}}
\def\lamp{\Blam_{\rho}}
\def\lamk{\Blam_{k}}
\def\lams{\Blam_{\sigma}}

\def\ba{{\bf b}_{\alpha}}
\def\bb{{\bf b}_{\beta}}
\def\fp{{\bf f}_{\rho}}
\def\fa{{\bf f}_{\alpha}}
\def\qa{{\bf q}_{\alpha}}
\def\qb{{\bf q}_{\beta}}
\def\za{{\bf z}_{\alpha}}
\def\zb{{\bf z}_{\beta}}
\def\wa{{\bf w}_{\alpha}}
\def\wb{{\bf w}_{\beta}}

\def\Na{\bar{\bf N}_{\alpha}}
\def\Nb{\bar{\bf N}_{\beta}}

\def\Up{{\bf U}_p}
\def\Un{{\bf U}_n}

\def\dFdl{\frac{\partial F}{\partial l}}
\def\dFddl{\frac{\partial F}{\partial \dot l}}

\def\Sr{s_\theta}
\def\Cr{c_\theta}
\def\Sp{s_\phi}
\def\Cp{c_\phi}
\def\Sy{s_\psi}
\def\Cy{c_\psi}
\def\Sa{s_{\alpha}}
\def\Ca{c_{\alpha}}
\def\Vp{v_{\phi}}


\iflatexml
\else
\usepackage{biblatex}
\addbibresource{references.bib}
\fi

\setcounter{tocdepth}{5}
\setcounter{secnumdepth}{3}

\title{\doctitle}
\ifdefined\maindoc
\author{John Lloyd and Antonio S\'anchez}
\setpubdate{Last update: March, 2022}

\iflatexml
\date{}
\fi
\fi

% graphics paths
\graphicspath{{./}{images/}}

% Listings settings
\definecolor{myblue}{rgb}{0,0,0.6}
\definecolor{mygreen}{rgb}{0,0.6,0}
\definecolor{mygray}{rgb}{0.5,0.5,0.5}
\definecolor{mylightgray}{rgb}{0.95,0.95,0.95}
\definecolor{mymauve}{rgb}{0.58,0,0.82}
\definecolor{myblack}{rgb}{0,0,0}
\lstset{
   language=Java,                   % text highlighting for Java
   breakatwhitespace=false,         % automatic breaks only at whitespace
   breaklines=true,                 % automatic line breaking
   commentstyle=\color{mygreen},    % comment style
   keepspaces=true,                 % keeps spaces in text
   keywordstyle=\color{myblue},     % keyword style
   numbers=none,                    % line-numbers; values: (none, left, right)
   numbersep=5pt,                   % how far the line-numbers are from code
   numberstyle=\tiny\color{mygray}, % line-numbers style
   showspaces=false,                % show spaces everywhere
   showstringspaces=false,          % underline spaces within strings
   showtabs=false,                  % show tabs
   stepnumber=1,                    % the step between two line-numbers
   stringstyle=\color{mymauve},     % string literal style
   tabsize=3,                       % sets default tabsize to 3 spaces
   backgroundcolor=\color{mylightgray}, % background color
   frame=single, 					% adds a frame around the code
   rulesepcolor=\color{mygray},
   rulecolor=\color{myblack},
   framerule=0pt,
   xleftmargin=2.2ex,               % numbers inside box
   framexleftmargin=2.2ex,			% indentation of frame
}

\begin{document}

\frontmatter

%\layout
\maketitle

\iflatexml{\large\pubdate}\fi

\tableofcontents

\mainmatter
\fi

\chapter{ArtiSynth Overview}
\label{Overview:sec}

ArtiSynth is an open-source, Java-based system for creating and
simulating mechanical and biomechanical models, with specific
capabilities for the combined simulation of rigid and deformable
bodies, together with contact and constraints. It is presently
directed at application domains in biomechanics, medicine, physiology,
and dentistry, but it can also be applied to other areas such as
traditional mechanical simulation, ergonomic design, and graphical and
visual effects.

\section{System structure}

An ArtiSynth model is composed of a hierarchy of models and model
components which are implemented by
various Java classes. These may include sub-models (including finite
element models), particles, rigid bodies, springs, connectors, and
constraints.  The component hierarchy may be in turn connected to
various {\it agent} components, such as control panels, controllers
and monitors, and input and output data streams (i.e., {\it probes}),
which have the ability to control and record the simulation as it
advances in time. Agents are presented in more detail in Section
\ref{SimulationControl:sec}. 

The models and agents are collected together within a top-level
component known as a {\it root model}.  Simulation proceeds under the
control of a {\it scheduler}, which advances the models through time
using a physics simulator. A rich graphical user interface (GUI)
allows users to view and edit the model hierarchy, modify component
properties, and edit and temporally arrange the input and output
probes using a {\it timeline} display.

\subsection{Model components}

Every ArtiSynth component is an instance of
\javaclass[artisynth.core.modelbase]{ModelComponent}. When connected
to the hierarchy, it is assigned a unique number relative to its
parent; the parent and number can be obtained using the methods
\javamethod[artisynth.core.modelbase.ModelComponent]{getParent()} and
\javamethod[artisynth.core.modelbase.ModelComponent]{getNumber()},
respectively.  Components may also be assigned a name (using
\javamethod*[artisynth.core.modelbase.ModelComponent]{setName()})
which is then returned using
\javamethod*[artisynth.core.modelbase.ModelComponent]{getName()}.

A sub-interface of {\tt ModelComponent} includes
\javaclass[artisynth.core.modelbase]{CompositeComponent}, which
contains child components.  A
\javaclass[artisynth.core.modelbase]{ComponentList} is a {\tt
CompositeComponent} which simply contains a list of other components
(such as particles, rigid bodies, sub-models, etc.).

Components which contain state information (such as position and
velocity) should extend 
\javaclass[artisynth.core.modelbase]{HasState}, 
which provides the methods
\javamethod*[artisynth.core.modelbase.HasState]{getState()}
and
\javamethod*[artisynth.core.modelbase.HasState]{setState()}
for saving and restoring state.

A
\javaclass[artisynth.core.modelbase]{Model}
is a sub-interface of {\tt CompositeComponent} and {\tt
HasState} that contains the notion of advancing through time and which
implements this with the methods {\tt initialize(t0)} and {\tt
advance(t0, t1, flags)}, as discussed further in
Section \ref{ModelAdvancement:sec}.
The most common instance of {\tt Model} used
in ArtiSynth is {\tt MechModel} (Section \ref{MechModel:sec}), which
is the top-level container for a mechanical or biomechanical model.

\subsection{The RootModel}
\label{RootModel:sec}

The top-level component in the hierarchy is the {\it root model},
which is a subclass of \javaclass[artisynth.core.workspace]{RootModel}
and which contains a list of models along with lists of agents used to
control and interact with these models. The component lists in {\tt
RootModel} include:

\begin{shadedregion}
\begin{tabular}{ll}
\tt models & top-level models of the component hierarchy \\
\tt inputProbes & input data streams for controlling the simulation \\
\tt controllers & functions for controlling the simulation \\
\tt monitors & functions for observing the simulation \\
\tt outputProbes & output data streams for observing the simulation \\
\end{tabular}
\end{shadedregion}

Each agent may be associated with a specific top-level model.

\subsection{Component path names}
\label{PathNames:sec}

The names and/or numbers of a component and its ancestors can be used to
form a component path name. This path has a construction 
analogous to Unix file path names, with the '/' character acting as a
separator. Absolute paths start with '/', which indicates the root
model. Relative paths omit the leading '/' and can begin lower down
in the hierarchy. A typical path name might be
\begin{verbatim}
  /models/JawHyoidModel/axialSprings/lad
\end{verbatim}
For nameless components in the path, their numbers can be used
instead.  Numbers can also be used for components that have
names. Hence the path above could also be represented using
only numbers, as in
\begin{verbatim}
  /0/0/1/5
\end{verbatim}
although this would most likely appear only in machine-generated
output.

\subsection{Model advancement}
\label{ModelAdvancement:sec}

ArtiSynth simulation proceeds by advancing all of the root model's
top-level models through a sequence of time steps. Every time
step is achieved by calling each model's 
\javamethod*[artisynth.core.modelbase.Model]{advance()} method:
\begin{lstlisting}[]
  public StepAdjustment advance (double t0, double t1) {
     ... perform simulation ...
  }
\end{lstlisting}
This method advances the model from time {\tt t0} to time {\tt t1},
performing whatever physical simulation is required (see Section
\ref{PhysicsSimulation:sec}). The method may optionally return a {\tt
StepAdjustment} indicating that the step size ({\tt t1} - {\tt t0}) was
too large and that the advance should be redone with a smaller step
size. 

The root model has it's own
\javamethod*[artisynth.core.workspace.RootModel]{advance()}, which in
turn calls the advance method for all of the top-level models, in
sequence. The advance of each model is surrounded by the application
of whatever 
agents are associated with that model. This is done
by calling the agent's {\tt apply()} method:
\begin{lstlisting}[]
   model.preadvance (t0, t1);
   for (each input probe p) {
      p.apply (t1);
   }
   for (each controller c) {
      c.apply (t0, t1);
   }
   model.advance (t0, t1);
   for (each monitor m) {
      m.apply (t0, t1);
   }
   for (each output probe p) {
      p.apply (t1);
   }
\end{lstlisting}
Agents not associated with a specific model are applied before (or
after) the advance of all other models.

More precise details about model advancement are given in the 
\artisynthManual{artisynth}{ArtiSynth Reference Manual}.

\subsection{MechModel}
\label{MechModel:sec}

Most ArtiSynth applications contain a single top-level model which is
an instance of \javaclass[artisynth.core.mechmodels]{MechModel}.  This
is a\pdfbreak
\javaclass[artisynth.core.modelbase]{CompositeComponent} that may
(recursively) contain an arbitrary number of mechanical components,
including finite element models, other {\tt MechModel}s, particles,
rigid bodies, constraints, attachments, and various force effectors.
The {\tt MechModel} {\tt advance()} method invokes a physics simulator
that advances these components forward in time (Section
\ref{PhysicsSimulation:sec}).

For convenience each {\tt MechModel} contains a number of predefined
containers for different component types, including:

%\begin{framed}%
%\colorbox{shadecolor}{%
\begin{shadedregion}
\begin{tabular}{ll}
\tt particles & 3 DOF particles \\
\tt points & other 3 DOF points \\
\tt rigidBodies & 6 DOF rigid bodies \\
\tt frames & other 6 DOF frames \\
\tt axialSprings & point-to-point springs \\
\tt connectors & joint-type connectors between bodies \\
\tt constrainers & general constraints \\
\tt forceEffectors & general force-effectors \\
\tt attachments & attachments between dynamic components \\
\tt renderables & renderable components (for visualization only) \\
\end{tabular}
\end{shadedregion}
Each of these is a child component of {\tt MechModel} and is
implemented as a
\javaclass[artisynth.core.modelbase]{ComponentList}. Special methods
are provided for adding and removing items from them. However,
applications are not required to use these containers, and may instead
create any component containment structure that is appropriate.
%(Section \ref{GeneralArrangements:sec}).  
If not used, the containers will simply remain empty.

\section{Physics simulation}
\label{PhysicsSimulation:sec}

Only a brief summary of ArtiSynth physics simulation is described
here.  Full details are given in \cite{lloyd2012artisynth} and in the
related
\href{http://www.artisynth.org/doc/artisynth.pdf}{overview paper}.

For purposes of physics simulation, the components of a {\tt
MechModel} are grouped as follows:

\begin{description}
	
\item[Dynamic components] \mbox{}\hfill\\ Components, such as a
particles and rigid bodies, that contain position and velocity state,
as well as mass. All dynamic components are instances of
the Java interface
\javaclass[artisynth.core.mechmodels]{DynamicComponent}.

\item[Force effectors] \mbox{}\hfill\\
Components, such as springs or finite elements,
that exert forces between dynamic components.
All force effectors are instances of the Java interface
\javaclass[artisynth.core.mechmodels]{ForceEffector}.

\item[Constrainers] \mbox{}\hfill\\
Components that enforce constraints between dynamic components. 
All constrainers are instances of the Java interface
\javaclass[artisynth.core.mechmodels]{Constrainer}.

\item[Attachments] \mbox{}\hfill\\
Attachments between dynamic components. While technically these
are constraints, they are implemented using a different approach.
All attachment components are instances of
\javaclass[artisynth.core.mechmodels]{DynamicAttachment}.

\end{description}

The positions, velocities, and forces associated with all the
dynamic components are denoted by the composite vectors 
$\q$, $\u$, and $\f$. 
In addition, the composite mass matrix is given by
$\M$. 
Newton's second law then gives
\begin{equation}
\f = \frac{d \M \u}{d t} = \M \dot\u + \dot\M \u,
\label{newton:eqn}
\end{equation}
%
where the $\dot\M \u$ accounts for various ``fictitious'' forces.

Each integration step involves solving for
the velocities $\u^{k+1}$ at time step $k+1$ given the velocities and forces
at step $k$. One way to do this is to solve the expression
%
\begin{equation}
\M \, \u^{k+1} = \M \u^k + h \bar\f
\label{velupdate:eqn}
\end{equation}
%
for $\u^{k+1}$, where $h$ is the step size and 
$\bar\f \equiv \f - \dot\M \u$. Given the updated velocities $\u^{k+1}$, one can
determine $\dot\q^{k+1}$ from
%
\begin{equation}
\dot\q^{k+1} = \Q \u^{k+1},
\end{equation}
%
where $\Q$ accounts for situations (like rigid bodies) where $\dot\q \ne
\u$, and then solve for the updated positions using 
%
\begin{equation}
\q^{k+1} = \q^k + h \dot \q^{k+1}.
\label{posupdate:eqn}
\end{equation}
%
(\ref{velupdate:eqn}) and (\ref{posupdate:eqn}) together comprise a
simple {\it symplectic Euler} integrator.

In addition to forces, bilateral and unilateral constraints give rise to
locally linear constraints on $\u$ of the form
%
\begin{equation}
\G(\q) \u = 0, \qquad \N(\q) \u \ge 0.
\label{constraints:eqn}
\end{equation}
%
Bilateral constraints may include rigid body joints, FEM
incompressibility, and point-surface constraints, while unilateral
constraints include contact and joint limits.  Constraints give rise
to constraint forces (in the directions $\G(\q)^T$ and $\N(\q)^T$)
which supplement the forces of (\ref{newton:eqn}) in order to enforce
the constraint conditions.  In addition, for unilateral constraints,
we have a complementarity condition in which $\N \u > 0$ implies no
constraint force, and a constraint force implies $\N \u = 0$.  Any
given constraint usually involves only a few dynamic components and so
$\G$ and $\N$ are generally sparse.

Adding constraints to the velocity solve (\ref{velupdate:eqn})
leads to a mixed linear complementarity problem (MLCP)
of the form
\begin{gather}
\left(
\begin{matrix}
\hat\M^{k} & -\G^{T} & -\N^{T} \\
\G & 0 & 0 \\
\N & 0 & 0 
\end{matrix}
\right)
\left(
\begin{matrix}
\u^{k+1} \\
\tilde\Blam \\
\tilde\Bthe
\end{matrix}
\right)
+
\left(
\begin{matrix}
-\M \u^{k} - h \hat\f^{k} \\
-\g \\
-\n
\end{matrix}
\right)
=
\left(
\begin{matrix}
0 \\
0 \\
\w
\end{matrix}
\right), \notag \\
0 \le \Bthe \perp \w \ge 0,
\label{KKTvelocity:eqn}
\end{gather}
where $\w$ is a slack variable, $\tilde\Blam$ and $\tilde\Bthe$ give the force
constraint impulses over the time step, and $\g$ and $\n$ are
derivative terms defined by
%
\begin{equation}
\g \equiv -h \dot\G \u^k, \quad \n \equiv -h \dot\N \u^k,
\end{equation}
%
to account for time variations in $\G$ and $\N$.
In addition,
$\hat\M$ and $\hat\f$ are $\M$ and $\bar\f$ augmented with stiffness
and damping terms terms to accommodate implicit integration, which
is often required for problems involving deformable bodies.
The actual constraint forces $\Blam$ and $\Bthe$ can be determined
by dividing the impulses by the time step $h$:
%
\begin{equation}
\Blam = \tilde\Blam/h, \quad \Bthe = \tilde\Bthe/h.
\label{impulsesToForces:eqn}
\end{equation}
%

We note here that ArtiSynth uses a {\it full coordinate} formulation,
in which the position of each dynamic body is solved using full, or
unconstrained, coordinates, with constraint relationships acting to
restrict these coordinates. In contrast, some other simulation
systems, including OpenSim \cite{delp2007opensim}, use {\it reduced}
coordinates, in which the system dynamics are formulated using a
smaller set of coordinates (such as joint angles) that implicitly take
the system's constraints into account. Each methodology has its own
advantages. Reduced formulations yield systems with fewer degrees of
freedom and no constraint errors. On the other hand, full coordinates
make it easier to combine and connect a wide range of components,
including rigid bodies and FEM models.

Attachments between components can be implemented by constraining the
velocities of the attached components using special constraints of the
form
%
\begin{equation}
\u_j = -\G_{j\alpha} \u_\alpha
\label{attachments:eqn}
\end{equation}
%
where $\u_j$ and $\u_\alpha$ denote the velocities of the attached and
non-attached components. The constraint matrix $\G_{j\alpha}$ is
sparse, with a non-zero block entry for each {\it master} component to
which the attached component is connected. The simplest case involves
attaching a point $j$ to another point $k$, with the simple velocity relationship
%
\begin{equation}
\u_j = \u_k
\end{equation}
%
That means that $\G_{j\alpha}$ has a single entry of $-\I$ (where $\I$
is the $3 \times 3$ identity matrix) in the $k$-th block column.
Another common case involves connecting a point $j$ to
a rigid frame $k$. The velocity relationship for this is
%
\begin{equation}
\u_j = \u_k - \l_j \times \Bom_k
\end{equation}
%
where $\u_k$ and $\Bom_k$ are the translational and rotational
velocity of the frame and $l_j$ is the location of the point relative
to the frame's origin (as seen in world coordinates). The corresponding
$\G_{j\alpha}$ contains a single $3 \times 6$ block entry of the form
%
\begin{equation}
\left(\begin{matrix}
\I & [ l_j ]
\end{matrix}\right)
\end{equation}
%
in the $k-th$ block column, where
%
\begin{equation}
[ l ] \equiv 
\left(\begin{matrix}
0 & -l_z & l_y \\
l_z & 0 & -l_x \\
-l_y & l_x & 0
\end{matrix}\right)
\end{equation}
%
is a skew-symmetric {\it cross product matrix}.
The attachment constraints $\G_{j\alpha}$ 
could be added directly to
(\ref{KKTvelocity:eqn}), but their special form allows us to
explicitly solve for $\u_j$, and hence reduce the size of
(\ref{KKTvelocity:eqn}), by factoring out the attached velocities
before solution.

The MLCP (\ref{KKTvelocity:eqn}) corresponds to a single step
integrator. However, higher order integrators, such as Newmark
methods, usually give rise to MLCPs with an equivalent form.  Most
ArtiSynth integrators use some variation of (\ref{KKTvelocity:eqn}) to
determine the system velocity at each time step.

To set up (\ref{KKTvelocity:eqn}), the {\tt MechModel} component
hierarchy is traversed and the methods of the different component
types are queried for the required values.  Dynamic components (type
{\tt DynamicComponent}) provide $\q$, $\u$, and $\M$; force effectors
({\tt ForceEffector}) determine $\hat\f$ and the stiffness/damping
augmentation used to produce $\hat\M$; constrainers ({\tt
Constrainer}) supply $\G$, $\N$, $\g$ and $\n$, and attachments ({\tt
DynamicAttachment}) provide the information needed to factor out
attached velocities.

\section{Basic packages}

The core code of the ArtiSynth project is divided into three main
packages, each with a number of sub-packages.

\subsection{maspack}

The packages under {\tt maspack} contain general computational
utilities that are independent of ArtiSynth and could be
used in a variety of other contexts. The main packages are:

\begin{lstlisting}[]
maspack.util               // general utilities
maspack.matrix             // matrix and linear algebra
maspack.graph              // graph algorithms
maspack.fileutil           // remote file access 
maspack.properties         // property implementation
maspack.spatialmotion      // 3D spatial motion and dynamics
maspack.solvers            // LCP solvers and linear solver interfaces
maspack.render             // viewer and rendering classes
maspack.geometry           // 3D geometry and meshes
maspack.collision          // collision detection
maspack.widgets            // Java swing widgets for maspack data types 
maspack.apps               // stand-alone programs based only on maspack
\end{lstlisting}

\subsection{artisynth.core}

The packages under {\tt artisynth.core} contain the core code for
ArtiSynth model components and its GUI infrastructure.

\begin{lstlisting}[]
artisynth.core.util        // general ArtiSynth utilities
artisynth.core.modelbase   // base classes for model components
artisynth.core.materials   // materials for springs and finite elements
artisynth.core.mechmodels  // basic mechanical models
artisynth.core.femmodels   // finite element models
artisynth.core.probes      // input and output probes
artisynth.core.workspace   // RootModel and associated components
artisynth.core.driver      // start ArtiSynth and drive the simulation
artisynth.core.gui         // graphical interface
artisynth.core.inverse     // inverse tracking controller
artisynth.core.moviemaker  // used for making movies
artisynth.core.renderables // components that are strictly visual
artisynth.core.opensim     // OpenSim model parser (under development)
artisynth.core.mfreemodels // mesh free models (experimental, not supported)
\end{lstlisting}

\subsection{artisynth.demos}

These packages contain demonstration models that illustrate
ArtiSynth's modeling capabilities:

\begin{lstlisting}[]
artisynth.demos.mech       // mechanical model demos
artisynth.demos.fem        // demos involving finite elements
artisynth.demos.inverse    // demos involving inverse control
artisynth.demos.tutorial   // demos in this manual
\end{lstlisting}

\section{Properties}
\label{Properties:sec}

ArtiSynth components expose {\it properties}, which provide a uniform
interface for accessing their internal parameters and
state. Properties vary from component to component; those for {\tt
RigidBody} include {\tt position}, {\tt orientation}, {\tt mass}, and
{\tt density}, while those for {\tt AxialSpring} include {\tt
restLength} and {\tt material}. Properties are particularly
useful for automatically creating control panels and
probes, as described in Section \ref{SimulationControl:sec}.
They are also used for automating component serialization.

Properties are described only briefly in this section; 
more detailed descriptions are available in the\pdfbreak
\artisynthManual{maspack}{Maspack Reference Manual} and the 
\href{http://www.artisynth.org/doc/artisynth.pdf}{overview paper}.

The set of properties defined for a component is fixed for that
component's class; while property values may vary between component
instances, their definitions are class-specific.  
Properties are exported by a class through code contained in the class
definition, as described in Section \ref{CustomProperties:sec}.

\subsection{Querying and setting property values}
\label{ControllingPropertyValues:sec}

Each property has a unique name that can be used to access its value
interactively in the GUI. This can be done either by using a custom
control panel (Section \ref{ControlPanels:sec}) or by selecting the
component and choosing {\sf Edit properties ...} from the right-click
context menu).

Properties can also be accessed in code using their set/get accessor
methods. Unless otherwise specified, the names for these are formed by
simply prepending {\tt set} or {\tt get} to the property's name.
More specifically, a property with the name {\sf foo} and a value type
of {\tt Bar} will usually have accessor signatures of
%
\begin{lstlisting}[]
  Bar getFoo()

  void setFoo (Bar value)
\end{lstlisting}
%

\subsection{Property handles and paths}
\label{PropertyHandlesAndPaths:sec}

A property's name can also be used to obtain a {\it
property handle} through which its value may be queried or
set generically. Property handles are implemented by
the class \javaclass[maspack.properties]{Property} and are returned by
the component's
\javamethod*[maspack.properties.HasProperties]{getProperty()} method.
{\tt getProperty()} takes a property's name and returns the
corresponding handle. For example, components of type {\tt Muscle}
have a property {\tt excitation}, for which a handle
may be obtained using a code fragment such as
\begin{lstlisting}[]
  Muscle muscle; 
  ...
  Property prop = muscle.getProperty ("excitation");
\end{lstlisting}
Property handles can also be obtained for
sub-components, using a {\it property path} that consists
of a path to the sub-component followed by a colon
`{\tt :}' and the property name. For example,
to obtain the {\tt excitation} property for a sub-component
located by {\tt axialSprings/lad} relative to a {\tt MechModel},
one could use a call of the form
\begin{lstlisting}[]
  MechModel mech;
  ...
  Property prop = mech.getProperty ("axialSprings/lad:excitation");
\end{lstlisting}

\subsection{Composite and inheritable properties}
\label{CompositeInheritableProperties:sec}

\begin{figure}[t]
\begin{center}
 \includegraphics[width=3.75in]{images/inheritedProperties}
\end{center}
\caption{Inheritance of a property named {\it stiffness} among
a component hierarchy. Explicit settings are in bold; inherited settings
are in gray italic.}
\label{inheritedProperties:fig}
\end{figure}

Composite properties are possible, in which a property value is a
composite object that in turn has sub-properties. A good example of
this is the {\tt RenderProps} class, which is
associated with the property {\tt renderProps} for renderable objects
and which itself can have a number of sub-properties such as {\tt
visible}, {\tt faceStyle}, {\tt faceColor}, {\tt lineStyle}, {\tt
lineColor}, etc.

Properties can be declared to be {\tt inheritable}, so that their
values can be inherited from the same properties hosted by ancestor
components further up the component hierarchy. Inheritable properties
require a more elaborate declaration and are associated with a {\it
mode} which may be either {\tt Explicit} or {\tt Inherited}.  If a
property's mode is inherited, then its value is obtained from
the closest ancestor exposing the same property whose mode is
explicit. In Figure (\ref{inheritedProperties:fig}), the property {\it
stiffness} is explicitly set in components A, C, and E, and inherited
in B and D (which inherit from A) and F (which inherits from C).

\section{Creating an application model}
\label{CreatingAnApplication:sec}

ArtiSynth applications are created by writing and compiling
an {\it application model} that is a subclass of {\tt RootModel}.
This application-specific root model is then loaded and run by the
ArtiSynth program.

The code for the application model should:

\begin{itemize}

\item Declare a no-args constructor

\item Override the {\tt RootModel}
\javamethod*[artisynth.core.workspace.RootModel]{build()}
method to construct the application.

\end{itemize}

ArtiSynth can load a model either using the build method
or by reading it from a file:

\begin{description}

\item[Build method] \mbox{}

ArtiSynth creates an instance of the
model using the no-args constructor, assigns it a name
(which is either user-specified or the simple name of the class), and
then calls the {\tt build()} method to perform the actual
construction.

\item[Reading from a file] \mbox{}

ArtiSynth creates an instance of the
model using the no-args constructor, and then the model is named
and constructed by reading the file.

\end{description}

The no-args constructor should perform whatever initialization is
required in both cases, while the {\tt build()} method takes the place of the
file specification. Unless a model is originally created using a file
specification (which is very tedious), the first time
creation of a model will almost always entail using the {\tt build()}
method.

The general template for application model code looks like this:

\begin{lstlisting}[]
package artisynth.models.experimental; // package where the model resides
import artisynth.core.workspace.RootModel;
... other imports ...

public class MyModel extends RootModel {

   // no-args constructor
   public MyModel() {
      ... basic initialization ...
   }

   // build method to do model construction
   public void build (String[] args) {
      ... code to build the model ....
   }
}
\end{lstlisting}
Here, the model itself is called {\tt MyModel}, and is defined in the
(hypothetical) 
package {\tt artisynth.models.experimental} (placing models in the super
package {\tt artisynth.models} is common practice but not
necessary).

\begin{sideblock}
Note: The {\tt build()} method was only introduced in ArtiSynth
3.1. Prior to that, application models were constructed using a
constructor taking a {\tt String} argument supplying the name of the
model. This method of model construction still works but is
deprecated.
\end{sideblock}

\subsection{Implementing the build() method}

As mentioned above, the {\tt build()} method is responsible for actual
model construction.  Many applications are built using a single
top-level {\tt MechModel}.  Build methods for these may look
like the following:
\begin{lstlisting}[]
   public void build (String[] args) {
      MechModel mech = new MechModel("mech");
      addModel (mech);

      ... create and add components to the mech model ...
      ... create and add any needed agents to the root model ...

   }
\end{lstlisting}
First, a \javaclass[artisynth.core.mechmodels]{MechModel} is created
(with the name {\tt "mech"} in this example, although any name, or no
name, may be given) and added to the list of models in the root model
using the {\tt addModel()} method. Subsequent code then creates and
adds the components required by the {\tt MechModel}, as described in
Sections \ref{MechModelsI:sec}, \ref{MechModelsII:sec} and
\ref{FEMModels:sec}.  The {\tt build()} method also creates and adds
to the root model any agents required by the application (controllers,
probes, etc.), as described in Section \ref{SimulationControl:sec}.

When constructing a model, there is no fixed order in which components
need to be added. For instance, in the above example, {\tt
addModel(mech)} could be called near the end of the {\tt build()}
method rather than at the beginning. The only restriction is that when
a component is added to the hierarchy, all other components that it
refers to should already have been added to the hierarchy. For
instance, an axial spring (Section \ref{ParticlesAndSprings:sec})
refers to two points. When it is added to the hierarchy, those two
points should already be present in the hierarchy.

The {\tt build()} method supplies a {\tt String} array as an argument,
which can be used to transmit application arguments in a manner
analogous to the {\tt args} argument passed to static {\tt main()}
methods. Build arguments can be specified when a model is loaded
directly from a class using {\sf Models > Load from class ...}, or when
the {\it startup model} is set to automatically load a model
when ArtiSynth is first started ({\sf Settings > Startup
model}). Details are given in the ``Loading, Simulating and
Saving Models'' section of the
\artisynthManual{uiguide}{User Interface Guide}.

Build arguments can also be listed directly on the ArtiSynth command
line when specifying a model to load using the {\tt -model
<classname>} option.  This is done by enclosing the desired arguments
within square brackets {\tt [ ]} immediately following the {\tt
-model} option. So, for example,
%
\begin{verbatim}
> artisynth -model projects.MyModel [ -size 50 ]
\end{verbatim}
%
would cause the strings {\tt "-size"} and {\tt "50"} to
be passed to the {\tt build()} method of {\tt MyModel}.

\subsection{Making models visible to ArtiSynth}

In order to load an application model into ArtiSynth, the classes
associated with its implementation must be made visible to ArtiSynth.
This usually involves adding the top-level class folder associated
with the application code to the classpath used by ArtiSynth.

\begin{sideblock}
The demonstration models referred to in this guide belong to the
package {\tt artisynth.demos.tutorial} and are already visible to
ArtiSynth.
\end{sideblock}

In most current ArtiSynth projects, classes are stored in
a folder tree separate from the source code, with the top-level
class folder named {\tt classes}, located one level below
the project root folder. A typical top-level class folder
might be stored in a location like this:
\begin{verbatim}
  /home/joeuser/artisynthProjects/classes
\end{verbatim}
In the example shown in Section \ref{CreatingAnApplication:sec}, the
model was created in the package {\tt artisynth.models.experimental}.
Since Java classes are arranged in a folder structure that mirrors
package names, with respect to the sample project folder shown
above, the model class would be located in
\begin{verbatim}
  /home/joeuser/artisynthProjects/classes/artisynth/models/experimental
\end{verbatim}

At present there are three ways to make top-level class folders
known to ArtiSynth:

\begin{description}

\item[Add projects to your Eclipse launch configuration]\mbox{}

If you are using the Eclipse IDE, then you can add the project in
which are developing your model code to the launch configuration that
you use to run ArtiSynth. Other IDEs will presumably provide similar
functionality.

\item[Add the folders to the external classpath]\mbox{}

You can explicitly add the class folders to ArtiSynth's external
classpath. The easiest way to do this is to select ``{\sf Settings >
External classpath ...}'' from the {\sf Settings} menu, which will
open an external classpath editor which lists all the classpath
entries in a large panel on the left. (When ArtiSynth is first
installed, the external classpath has no entries, and so this panel
will be blank.) Class folders can then by added via
the {\sf ``Add class folder''} button, and the classpath is
saved using the {\sf Save} button.

\item[Add the folders to your CLASSPATH environment variable]\mbox{}

If you are running ArtiSynth from the command line, using the {\tt
artisynth} command (or {\tt artisynth.bat} on Windows), then you can
define a CLASSPATH environment variable in your environment and add
the needed folders to this.

\end{description}

All of these methods are described in more detail in the ``Installing
External Models and Packages'' section of the ArtiSynth Installation
Guide (available for 
\artisynthManual[installation/linuxInstallation]{linuxInstallation}{Linux}, 
\artisynthManual[installation/windowsInstallation]{windowsInstallation}{Windows},
and 
\artisynthManual[installation/macosInstallation]{macosInstallation}{MacOS}).

\subsection{Loading and running a model}
\label{LoadingAndRunning:sec}

If a model's classes are visible to ArtiSynth, then it may be loaded
into ArtiSynth in several ways:

\begin{description}

\item[Loading from the Model menu]\mbox{}

If the root model is contained in a package located under {\tt
artisynth.demos} or {\tt artisynth.models}, then it will appear in the
default model menu ({\sf Models} in the main menu bar) under the
submenu {\sf All demos} or {\sf All models}.

\item[Loading by class path]\mbox{}

A model may also be loaded by choosing {\sf ``Load from class ...''}
from the {\sf Models} menu and specifying its package name and then
choosing its root model class. It is also possible to use the {\tt
-model <classname>} command line argument to have a model loaded
directly into ArtiSynth when it starts up.

\item[Loading from a file]\mbox{}

If a model has been saved to a {\tt .art} file, it may be loaded from
that file by choosing {\sf File > Load model ...}.

\end{description}

These methods are described in detail in the 
section ``Loading and Simulating Models'' of the\pdfbreak
\artisynthManual{uiguide}{ArtiSynth User Interface Guide}.

\begin{sideblock}
The demonstration models referred to in this guide should already
be present in the model menu and may be loaded
from the submenu {\sf Models > All demos > tutorial}.
\end{sideblock}

Once a model is loaded, it can be simulated, or {\it run}.  Simulation
of the model can then be started, paused, single-stepped, or reset
using the play controls (Figure \ref{PlayControlsFig}) located at the
upper right of the ArtiSynth window frame.  Starting and stopping a
simulation is done by clicking {\sf play/pause}, while {\sf reset}
resets the simulation to time 0.  The {\sf single-step} button
advances the simulation by one time step. The {\sf stop-all} button
will also stop the simulation, along with any Jython commands or
scripts that are running.

\begin{figure}[ht]
\begin{center}
\iflatexml
\includegraphics[]{../uiguide/images/playControls}
\else
\includegraphics[width=3in]{../uiguide/images/playControls}
\fi
\end{center}
\caption{The ArtiSynth play controls. From left to right: step size
control, current simulation time, and the reset, skip-back,
play/pause, single-step, skip-forward and stop-all buttons.}%
\label{PlayControlsFig}
\end{figure}

Comprehensive information on exploring and interacting with models is
given in the
\artisynthManual{uiguide}{ArtiSynth User Interface Guide}.

\ifdefined\maindoc
\else
\end{document}
\fi
